\documentclass[../../main.tex]{subfiles}
\begin{document}
\section{Riemann 距离函数}\label{section:ch13-distance-function}

Riemann 度量给予我们的最重要工具之一，是定义曲线长度的能力。设 $(M,g)$ 是一个带边或不带边的 Riemann 流形。若 $\gamma:[a,b]\to M$ 是一段分段光滑曲线，则 $\gamma$ 的长度为
\begin{align*}
L_g(\gamma)=\int_a^b |\gamma'(t)|_g\,\mathrm{d}t.
\end{align*}
因为 $|\gamma'(t)|_g$ 除有限多个 $t$ 值外都连续，并且在这些点处有良好定义的左右极限，所以该积分是良定义的。

\begin{exercise}\label{exercise:ch13-23}
设 $\gamma:[a,b]\to M$ 是一段分段光滑曲线，且 $a<c<b$。证明
\begin{align*}
L_g(\gamma)=L_g(\gamma|_{[a,c]})+L_g(\gamma|_{[c,b]}).
\end{align*}
\end{exercise}

\begin{exercise}\label{exercise:ch13-24}
证明曲线长度是 Riemann 流形的局部等距不变量。更确切地说，设 $(M,g)$ 与 $(\widetilde M,\widetilde g)$ 是带边或不带边的 Riemann 流形，$F:M\to\widetilde M$ 是一个局部等距。证明对 $M$ 中每一段分段光滑曲线 $\gamma$，都有 $L_{\widetilde g}(F\circ\gamma)=L_g(\gamma)$。
\end{exercise}

一个极其重要的事实是：长度在下述意义下与参数化无关。在第~11 章中，我们把一段分段光滑曲线 $\gamma:[a,b]\to M$ 的\emph{重新参数化}定义为形如 $\widetilde\gamma=\gamma\circ\varphi$ 的曲线段，其中 $\varphi:[c,d]\to[a,b]$ 是一个微分同胚。

\begin{proposition}[长度的参数无关性]\label{proposition:ch13-length-param-indep}
设 $(M,g)$ 是一个带边或不带边的 Riemann 流形，$\gamma:[a,b]\to M$ 是一段分段光滑曲线。若 $\widetilde\gamma$ 是 $\gamma$ 的一个重新参数化，则 $L_g(\widetilde\gamma)=L_g(\gamma)$。
\end{proposition}

\begin{proof}
首先设 $\gamma$ 光滑，$\varphi:[c,d]\to[a,b]$ 是一个使 $\widetilde\gamma=\gamma\circ\varphi$ 的微分同胚。$\varphi$ 是微分同胚这一事实意味着处处有 $\varphi'>0$ 或 $\varphi'<0$。先设 $\varphi'>0$。我们有
\begin{align*}
L_g(\widetilde\gamma)
&=\int_c^d |\widetilde\gamma'(t)|_g\,\mathrm{d}t
=\int_c^d \left|\frac{\mathrm{d}}{\mathrm{d}t}(\gamma\circ\varphi)(t)\right|_g\mathrm{d}t\\
&=\int_c^d |\varphi'(t)\gamma'(\varphi(t))|_g\,\mathrm{d}t
=\int_c^d |\gamma'(\varphi(t))|_g\varphi'(t)\,\mathrm{d}t\\
&=\int_a^b |\gamma'(s)|_g\,\mathrm{d}s=L_g(\gamma),
\end{align*}
其中倒数第二个等号用了积分的变量代换公式。在 $\varphi'<0$ 的情形，只需在上面的计算中引入两次符号改变即可。第一次是把 $\varphi'(t)$ 移出绝对值号时，因为 $|\varphi'(t)|=-\varphi'(t)$。第二次是在变量代换时，因为 $\varphi$ 反转了积分方向。由于两次符号改变相互抵消，结果相同。

若 $\gamma$ 只是分段光滑，则只需在它光滑的每个子区间上应用同样的论证。
\end{proof}

把曲线段当作“卷尺”，我们便能定义 Riemann 流形上点与点之间的距离。设 $(M,g)$ 是一个连通的 Riemann 流形。（当 $\partial M=\varnothing$ 时理论最直接，故本节余下部分作此假设。）对任意 $p,q\in M$，从 $p$ 到 $q$ 的（Riemann）\emph{距离}，记作 $d_g(p,q)$，定义为所有从 $p$ 到 $q$ 的分段光滑曲线段 $\gamma$ 的 $L_g(\gamma)$ 的下确界。由于 $M$ 中任意两点都可用一段分段光滑曲线连接（命题\ref{proposition:ch11-join-curves}），这是良定义的。

\begin{example}\label{example:ch13-euclidean-distance}
在 $(\mathbb{R}^n,\bar g)$ 中，习题\ref{exercise:ch13-10}表明任意直线段都是其端点之间最短的分段光滑曲线段。因此距离函数 $d_{\bar g}$ 等于通常的 Euclid 距离：
\begin{align*}
d_{\bar g}(x,y)=|x-y|.
\end{align*}
\end{example}

\begin{exercise}\label{exercise:ch13-27}
设 $(M,g)$ 与 $(\widetilde M,\widetilde g)$ 是连通的 Riemann 流形，$F:M\to\widetilde M$ 是一个 Riemann 等距。证明对所有 $p,q\in M$，都有 $d_{\widetilde g}(F(p),F(q))=d_g(p,q)$。
\end{exercise}

我们将在下面看到，Riemann 距离函数把 $M$ 变成一个度量空间，其拓扑与给定的流形拓扑相同。关键是下面这个技术性引理，它表明在坐标中每个 Riemann 度量都局部地与 Euclid 度量可比。

\begin{lemma}\label{lemma:ch13-local-comparability}
设 $g$ 是开子集 $U\subseteq\mathbb{R}^n$ 上的一个 Riemann 度量。给定紧子集 $K\subseteq U$，存在正常数 $c,C$，使得对所有 $x\in K$ 和所有 $v\in T_x\mathbb{R}^n$，
\begin{align}
c|v|_{\bar g}\leqslant |v|_g\leqslant C|v|_{\bar g}.\label{align:ch13-local-comparability}
\end{align}
\end{lemma}

\begin{proof}
对任意紧子集 $K\subseteq U$，令 $L\subseteq T\mathbb{R}^n$ 为集合
\begin{align*}
L=\{(x,v)\in T\mathbb{R}^n:x\in K,\ |v|_{\bar g}=1\}.
\end{align*}
在 $T\mathbb{R}^n$ 与 $\mathbb{R}^n\times\mathbb{R}^n$ 的典范等同下，$L$ 恰是乘积集 $K\times S^{n-1}$，因而是紧的。由于范数 $|v|_g$ 在 $L$ 上连续且严格为正，存在正常数 $c,C$，使当 $(x,v)\in L$ 时 $c\leqslant|v|_g\leqslant C$。若 $x\in K$，$v$ 是 $T_x\mathbb{R}^n$ 中的非零向量，令 $\lambda=|v|_{\bar g}$。则 $(x,\lambda^{-1}v)\in L$，由范数的齐次性，
\begin{align*}
|v|_g=\lambda\,|\lambda^{-1}v|_g\leqslant\lambda C=C|v|_{\bar g}.
\end{align*}
类似计算给出 $|v|_g\geqslant c|v|_{\bar g}$。当 $v=0$ 时同样的不等式显然成立。
\end{proof}

\begin{theorem}[作为度量空间的 Riemann 流形]\label{theorem:ch13-metric-space}
设 $(M,g)$ 是一个连通的 Riemann 流形。带上 Riemann 距离函数后，$M$ 是一个度量空间，其度量拓扑与原来的流形拓扑相同。
\end{theorem}

\begin{proof}
由定义立即有 $d_g(p,q)\geqslant0$。因为每个常值曲线段长度为零，故 $d_g(p,p)=0$；而 $d_g(p,q)=d_g(q,p)$ 来自这样一个事实：任何从 $p$ 到 $q$ 的曲线段都可重新参数化为从 $q$ 到 $p$。设 $\gamma_1,\gamma_2$ 分别是从 $p$ 到 $q$、从 $q$ 到 $r$ 的分段光滑曲线段（图\ref{figure:ch13-triangle}），令 $\gamma$ 是一段先沿 $\gamma_1$ 再沿 $\gamma_2$ 的分段光滑曲线段（必要时重新参数化）。则
\begin{align*}
d_g(p,r)\leqslant L_g(\gamma)=L_g(\gamma_1)+L_g(\gamma_2).
\end{align*}
对所有这样的 $\gamma_1,\gamma_2$ 取下确界，得 $d_g(p,r)\leqslant d_g(p,q)+d_g(q,r)$。（这就是为什么要用分段光滑曲线而不只是光滑曲线来定义距离函数很重要的原因之一。）

为完成 $(M,d_g)$ 是度量空间的证明，只需证明当 $p\ne q$ 时 $d_g(p,q)>0$。为此，设 $p,q\in M$ 为不同点，令 $U$ 是一个含 $p$ 但不含 $q$ 的光滑坐标域。如常利用坐标映射把 $U$ 等同于 $\mathbb{R}^n$ 的一个开子集，并令 $\bar g$ 表示这些坐标下的 Euclid 度量。若 $V$ 是以 $p$ 为心、半径为 $\varepsilon$ 的一个正则坐标球且 $\bar V\subseteq U$，引理\ref{lemma:ch13-local-comparability}表明存在正常数 $c,C$，使当 $x\in\bar V$、$v\in T_xM$ 时满足\eqref{align:ch13-local-comparability}。于是对任意完全含于 $\bar V$ 内的分段光滑曲线段 $\gamma$，有
\begin{align*}
cL_{\bar g}(\gamma)\leqslant L_g(\gamma)\leqslant CL_{\bar g}(\gamma).
\end{align*}
设 $\gamma:[a,b]\to M$ 是一段从 $p$ 到 $q$ 的分段光滑曲线。令 $t_0$ 为所有使 $\gamma(t)\notin\bar V$ 的 $t\in[a,b]$ 的下确界（图\ref{figure:ch13-positivity}）。由连续性，$\gamma(t_0)\in\partial V$，且当 $a\leqslant t\leqslant t_0$ 时 $\gamma(t)\in\bar V$。于是
\begin{align*}
L_g(\gamma)\geqslant L_g(\gamma|_{[a,t_0]})
\geqslant cL_{\bar g}(\gamma|_{[a,t_0]})
\geqslant c\,d_{\bar g}(p,\gamma(t_0))=c\varepsilon.
\end{align*}
对所有这样的 $\gamma$ 取下确界，得 $d_g(p,q)\geqslant c\varepsilon>0$。

\begin{figure}[htbp]
\centering
\begin{minipage}{0.48\textwidth}
\centering
\includegraphics[width=\linewidth]{fig13-2-triangle.png}
\caption{三角不等式}
\label{figure:ch13-triangle}
\end{minipage}
\end{figure}

\begin{figure}[htbp]
\centering
\begin{minipage}{0.48\textwidth}
\centering
\includegraphics[width=\linewidth]{fig13-3-positivity.png}
\caption{$d_g$ 的正定性}
\label{figure:ch13-positivity}
\end{minipage}
\end{figure}

最后，为证明 $d_g$ 生成的度量拓扑与 $M$ 上给定的流形拓扑相同，需证明流形拓扑中的开子集在度量拓扑中开，反之亦然。首先设 $U\subseteq M$ 在流形拓扑中开。令 $p\in U$，取如上以 $p$ 为心、半径 $\varepsilon$ 且 $\bar V\subseteq U$ 的正则坐标球 $V$。上一段论证表明，只要 $q\notin\bar V$，就有 $d_g(p,q)\geqslant c\varepsilon$。这一命题的逆否命题是：$d_g(p,q)<c\varepsilon$ 蕴含 $q\in\bar V\subseteq U$，换言之，以 $p$ 为心、半径 $c\varepsilon$ 的度量球含于 $U$。这表明 $U$ 在度量拓扑中开。

反之，设 $W$ 在度量拓扑中开，令 $p\in W$。取以 $p$ 为心、半径 $r$ 的正则坐标球 $V$，令 $\bar g$ 为给定坐标在 $\bar V$ 上确定的 Euclid 度量，令 $c,C$ 为使 $v\in T_qM$、$q\in\bar V$ 时满足\eqref{align:ch13-local-comparability}的正常数。取一个充分小的正数 $\varepsilon<r$，使以 $p$ 为心、半径 $C\varepsilon$ 的度量球含于 $W$；令 $V_\varepsilon$ 为 $\bar V$ 中与 $p$ 的 Euclid 距离小于 $\varepsilon$ 的点集。若 $q\in V_\varepsilon$，令 $\gamma$ 为坐标中从 $p$ 到 $q$ 的直线段。如上利用引理\ref{lemma:ch13-local-comparability}，得
\begin{align*}
d_g(p,q)\leqslant L_g(\gamma)\leqslant CL_{\bar g}(\gamma)<C\varepsilon.
\end{align*}
这表明 $V_\varepsilon$ 含于以 $p$ 为心、半径 $C\varepsilon$ 的度量球内，因而含于 $W$。由于 $V_\varepsilon$ 是 $p$ 在流形拓扑中的一个邻域，这表明 $W$ 在流形拓扑中也开。
\end{proof}

作为这个定理的推论，度量空间的全部术语都可以移植到连通 Riemann 流形上。于是，若 $(M,d_g)$ 是一个完备度量空间（即 $M$ 中每个 Cauchy 列都收敛到 $M$ 中一点），就说连通 Riemann 流形 $(M,g)$ 是\emph{完备的}，而 $g$ 是一个\emph{完备 Riemann 度量}；若存在常数 $K$ 使对所有 $x,y\in B$ 都有 $d_g(x,y)\leqslant K$，就说子集 $B\subseteq M$ 是\emph{有界的}。习题\ref{exercise:ch13-17}和习题\ref{exercise:ch13-18}概述了两个不同的证明，说明每个连通光滑流形都容许一个完备 Riemann 度量。

回忆一下，若一个拓扑空间容许一个距离函数，其度量拓扑与给定拓扑相同，就说它是\emph{可度量化的}。

\begin{corollary}\label{corollary:ch13-metrizable}
每个带边或不带边的光滑流形都是可度量化的。
\end{corollary}

\begin{proof}
首先设 $M$ 是一个不带边的光滑流形，任取 $M$ 上一个 Riemann 度量 $g$。若 $M$ 连通，定理\ref{theorem:ch13-metric-space}表明 $M$ 可度量化。更一般地，令 $\{M_i\}$ 为 $M$ 的连通分支，对每个 $i$ 取一点 $p_i\in M_i$。对 $x\in M_i$、$y\in M_j$，当 $i=j$ 时如定理\ref{theorem:ch13-metric-space}那样定义 $d_g(x,y)$，否则定义
\begin{align*}
d_g(x,y)=d_g(x,p_i)+1+d_g(p_j,y).
\end{align*}
（设想在不同分支中每对点 $p_i,p_j$ 之间架起一座长度为 $1$ 的“桥”，于是从 $x$ 到 $y$ 必须先到 $p_i$，过桥到 $p_j$，再从 $p_j$ 到 $y$。）容易验证这是一个在 $M$ 上诱导给定拓扑的距离函数。最后，若 $M$ 有非空边界，只需把 $M$ 嵌入它的倍流形（例\ref{example:ch09-double}），并注意可度量化拓扑空间的子空间总是可度量化的。
\end{proof}
\end{document}
